% ------------------------------------------------------------
\subsection{Interpretação dos Resultados}
\label{subsec:interpretacao_resultados}
% ------------------------------------------------------------

Os resultados confirmam o \textit{XGBoost} como o algoritmo mais adequado ao
rastreio do GAD nesta base, com o maior Kappa (0,372) e a maior AUC (0,751).
A vantagem do \textit{boosting} sobre o SVM e a ADTree evidencia que a
capacidade de modelar interações não lineares entre poucos sintomas pesa mais,
neste problema, do que a fronteira de decisão do hiperplano ou da árvore
alternada.

A escolha da técnica de balanceamento não define um vencedor absoluto, mas um
compromisso (\textit{trade-off}) atrelado ao objetivo clínico. O
\textit{BorderlineSMOTE} maximiza a concordância e preserva a Especificidade,
sendo preferível quando um falso positivo implica encaminhamento desnecessário.
Já o \textit{EasyEnsemble} eleva a Sensibilidade para 68,5\%, mostrando-se mais
adequado ao rastreio populacional em larga escala, no qual não detectar um caso
real é o erro mais grave.

O achado mais importante, contudo, é o teto de desempenho imposto pela amostra.
O experimento de Monte Carlo mostrou que os \textit{hard samples} permanecem
intrinsecamente difíceis mesmo sob reamostragem repetida, e a correção do
vazamento de dados (que derrubou o Kappa de 0,81 para $-0,08$) confirma que a
limitação vem da escassez de sinal na classe positiva, não da instabilidade do
treino. Como nenhuma configuração ultrapassou a concordância razoável
(\textit{fair}), os modelos devem ser lidos como ferramentas auxiliares de
triagem, capazes de priorizar crianças em risco, mas não de substituir a
avaliação clínica especializada.
